
achève par
récurrence, en notant que si $0\leq s\leq r$, et si $x_j=0$ pour $s<j\leq r$, on aura :
\[
  \sum_{i=0}^{s}f^*(x_i)\xi^i=0,
\]
d'où :
\[
  \sum_{i=0}^{s}f^*(x_i)\xi^{i+r-s}=0,
\]
d'où $x_s=0$ d'après ce qui précède.

\section*{3. Conséquences du théorème de structure pour le \(K^{\bullet}\) d'un fibré projectif}

\textbf{Proposition 3.1.} \uline{Soient} \(S\) \uline{un schéma quasi}-\uline{compact}, \(\underline{E}_1,\ldots,\underline{E}_n\) \uline{des} \(\underline{O}_S\)\uline{-Modules localement libres de type fini}, \(\mathbf D(\underline{E}_i)\) \uline{le schéma des drapeaux de} \(\underline{E}_i\), \(X=\mathbf D(\underline{E}_1)\times_S\ldots\times_S \mathbf D(\underline{E}_n)\).
\uline{Alors l'homomorphisme canonique}~:
\[
  K^{\bullet}(S)\longrightarrow K^{\bullet}(X)
\]
\uline{est injectif} (\(K^{\bullet}\) \uline{étant défini}, \uline{soit à partir des} \(\underline{O}_S\)\uline{-Modules localement libres de type fini}, \uline{soit à partir des complexes parfaits}).

On vérifie immédiatement en utilisant la propriété universelle des schémas de drapeaux que \(X\) peut être construit comme suit : on considère la suite de \(S\)-schémas \(X_1,\ldots,X_n\), définis par récurrence en prenant pour \(X_k\) le schéma de drapeaux de l'image inverse de \(\underline{E}_k\) sur \(X_{k-1}\); alors \(X=X_n\). On voit alors tout de suite que pour démontrer la proposition, il suffit de la montrer dans le cas où \(n=1\), le cas général s'en déduisant par récurrence. On peut donc supposer que \(X\) est le schéma des drapeaux d'un \(\underline{O}_S\)-Module localement libre de type fini \(\underline{E}\).

On sait que \(X\) peut alors être construit comme suit : on pose \(\underline{E}=\underline{F}_0\), \(S=X_0\), et on définit par récurrence des schémas \(X_1,\ldots,X_n\) et
